%=====================================================================
% PART II OPENER
%=====================================================================
\thispagestyle{plain}
% Part-opener maps get their own figure prefix (P2.1, ...).
\structuralfigures{P2.}

\begin{center}
\begin{tcolorbox}[enhanced, width=0.92\textwidth, colback=secondary!5,
  colframe=secondary, boxrule=0pt, leftrule=5pt, arc=3pt, top=8pt, bottom=8pt]
\large\itshape\color{secondary!80!black}
Most of a data scientist's working life is spent here: making messy data usable,
looking at it honestly, and being careful about what a sample is entitled to say
about a population. No model can rescue a project that skipped this part.
\end{tcolorbox}
\end{center}

\vspace{6pt}

\dsfigh{%
\begin{tikzpicture}[
  cbox/.style={rounded corners=3pt, draw=secondary, line width=1pt, fill=secondary!7,
               text=secondary!60!black, font=\small\bfseries, align=center,
               minimum width=3.4cm, minimum height=1.05cm},
  arr/.style={-{Stealth[length=2.4mm]}, secondary!60, line width=1pt},
  note/.style={font=\scriptsize\itshape, text=gray!60!black, align=center}
]
  \node[cbox] (c7)  at (0,0)     {Ch.\ 7\\Wrangling};
  \node[cbox] (c8)  at (4.3,0)   {Ch.\ 8\\Exploring};
  \node[cbox] (c9)  at (8.6,0)   {Ch.\ 9\\Inferring};
  \node[cbox] (c10) at (12.9,0)  {Ch.\ 10\\Explaining};

  \draw[arr] (c7) -- (c8);
  \draw[arr] (c8) -- (c9);
  \draw[arr] (c9) -- (c10);

  \node[note] at (0,-1.15)    {raw $\rightarrow$ usable};
  \node[note] at (4.3,-1.15)  {usable $\rightarrow$ understood};
  \node[note] at (8.6,-1.15)  {sample $\rightarrow$ population};
  \node[note] at (12.9,-1.15) {association $\rightarrow$ cause};
\end{tikzpicture}}%
{Reading path through Part II. Each chapter takes the data one step further from its raw
state towards a defensible claim.}{part2map}

\newpage
